\section{OMNeT++: architettura e workflow sperimentale}
\label{env:omnet}

OMNeT++ è il simulatore a eventi discreti che costituisce il nucleo dell’ambiente sperimentale adottato in questa tesi. La sua rilevanza nel contesto della ricerca sulle reti deriva dall’architettura modulare e dal supporto alla composizione gerarchica dei modelli, che consente di descrivere sistemi complessi come insiemi di moduli connessi tra loro tramite messaggi ed eventi temporizzati \cite{varga2008overview}. In ambito V2X, tale modularità è particolarmente utile perché permette di integrare stack di comunicazione, applicazioni, logiche di controllo e interfacce verso simulatori esterni di traffico.

Dal punto di vista architetturale, OMNeT++ separa:
\begin{itemize}
\item la descrizione strutturale del sistema (topologia di moduli e connessioni), tipicamente espressa tramite file NED;
\item la logica di comportamento dei moduli, implementata in C++;
\item la configurazione delle simulazioni e dei parametri sperimentali, gestita tramite file \texttt{omnetpp.ini}.
\end{itemize}
Questa separazione è particolarmente vantaggiosa per una tesi sperimentale, perché facilita la riproduzione di scenari varianti (parametri radio, policy di controllo, durata della simulazione, seed, logging) senza modificare continuamente il codice.

Nel workflow adottato, OMNeT++ svolge il ruolo di orchestratore dell’intera simulazione. Le principali attività gestite a questo livello sono:
\begin{itemize}
\item inizializzazione dei moduli e lettura della configurazione;
\item schedulazione degli eventi di rete e delle applicazioni;
\item raccolta di statistiche e segnali;
\item coordinamento con i moduli di mobilità e con le librerie esterne (Veins, PLEXE, Simu5G).
\end{itemize}
In altre parole, OMNeT++ non viene utilizzato come semplice “esecutore” di un singolo protocollo, ma come piattaforma su cui comporre un ambiente sperimentale multi-dominio.

Un ulteriore punto di forza, rilevante per il presente lavoro, è la possibilità di integrare framework sviluppati dalla comunità (ad es.\ Veins, PLEXE, Simu5G) mantenendo una struttura relativamente modulare. Questo aspetto è cruciale quando si devono aggiornare versioni dei componenti o sostituire un sottosistema (come la componente cellulare) senza riscrivere l’intera infrastruttura sperimentale. La letteratura più recente sul cooperative driving continua infatti a fare leva su questa modularità per costruire toolchain realistiche e riusabili \cite{segata2023multi-technology,nardini2020simu5g}.
